\newpage
\section*{Kurzfassung}
\vspace*{-0.8em}
\addcontentsline{toc}{section}{Kurzfassung}
\selectlanguage{ngerman}
Die fortschreitende Digitalisierung industrieller Produktion erfordert dynamische Planungsansätze, in denen Digitale Zwillinge als Bindeglied zwischen physischer Anlage und virtuellem Modell fungieren. Die vorliegende Arbeit untersucht die Potenziale und Herausforderungen, die sich aus der Nutzung von NVIDIA Omniverse sowie der darauf aufbauenden Simulationsplattform Isaac Sim (im Folgenden: Isaac Sim) für die Fabrikplanung ergeben.

Im Labor für Produktentwicklung und Digitale Fabrik der Hochschule Hannover werden auf Basis dieser Plattformen Digitale Zwillinge zweier realer Anlagen entwickelt: ein Rundtakttisch („Maze Runner“) sowie der Gelenkarmroboter „Dobot Magician“. Die Modellierung erfolgt unter Nutzung des offenen Datenformats OpenUSD sowie der Middleware-Architekturen OPC UA, MQTT und ROS 2 in Rahmen einer Python basierten Extension für NVIDIA Omniverse.

Ergänzend wird eine methodische Anleitung zur Reproduzierbarkeit der Modellerstellung entwickelt, sodass Studierende und industrielle Anwender den Prozess eigenständig nachvollziehen können. Die Arbeit schließt mit einer kritischen Diskussion infrastruktureller, ökonomischer und qualifikatorischer Anforderungen sowie der Formulierung strategischer Handlungsempfehlungen für Industrieunternehmen und Hochschulen.\\
\textbf{Schlagworte:} Digitaler Zwilling, NVIDIA Omniverse, Isaac Sim, OpenUSD, Fabrikplanung, virtuelle Inbetriebnahme, OPC UA, ROS\,2.

\vspace{0.6em}
\section*{Abstract}
\vspace*{-0.8em}
\addcontentsline{toc}{section}{Abstract}
\selectlanguage{english}
The ongoing digitalisation of industrial production calls for dynamic planning approaches in which digital twins serve as the interface between physical assets and their virtual representation. This thesis investigates the potentials and challenges that arise from applying NVIDIA Omniverse and its simulation platform Isaac Sim (hereafter referred to as Isaac Sim) to factory planning. Within the Laboratory of Product Development and Digital Factory at Hannover University of Applied Sciences and Arts, digital twins of two real installations are developed: a rotary indexing table (\enquote{Maze Runner}), provided by Yübo Wang (Hochschule RheinMain\,/\,Mercedes-Benz), who also contributed to shaping the contents of this project, and the articulated robot arm \enquote{Dobot Magician}. Modelling is based on the open data format OpenUSD and on the middleware architectures OPC UA, MQTT and ROS\,2, implemented within a Python-based extension for NVIDIA Omniverse. A methodological guideline ensures reproducibility, enabling students and industrial users to recreate the modelling workflow autonomously. The thesis concludes with a critical discussion of infrastructural, economic, and qualificational requirements as well as strategic recommendations for industry and higher education.

\textbf{Keywords:} digital twin, NVIDIA Omniverse, Isaac Sim, OpenUSD, factory planning, virtual commissioning, OPC UA, ROS\,2.\\
\selectlanguage{ngerman}
